\chapter{Related Work} \label{chapter:related-work}
This chapter provides an overview of the related research within the field of AI-assisted software development and robotic application development.

\section{Methodology}
This section describes the methodology used to identify and analyze the presented literature below. A systematic search was conducted across the SDU Library database, which include numerous scientific databases. \cite{SDU-Library} However, the search produced few directly relevant studies on the use of LLMs for development of robotic programs. Thus, the search was split into two focused systematic searches, one on AI-assisted software development and one on robotic solution development and reuse of data within robotics. The search was conducted using SDU Library and Google Scholar with keywords such as "AI assisted development" and "robotic data reuse". \cite{google-scholar} The search was supplemented with reverse snowballing. Only English language papers were included, and newer studies with more citations was prioritized. Studies were screened based on title, abstract, and full-text assessment for relevance.


\section{AI-Assisted Software Development}
Large language models have changed software development by making code generation and reasoning accessible through natural language interaction. AI-powered code assistants such as GitHub Copilot embed these capabilities directly in development environments, where they assist with development from local context. \cite{GitHubCopilot2026} Empirical studies report measurable productivity gains for routine and well-specified tasks, while benchmark evaluations show that code models can solve a meaningful subset of programming problems, although performance still depends strongly on task complexity and specification quality. \cite{bird2022taking, chen2021evaluating, austin2021program}

Their use extends beyond code synthesis. AI assistants are used for testing, debugging, maintenance, documentation, and explanation of unfamiliar code or technologies, which positions them as interactive knowledge aids as well as automation tools. \cite{barke2023grounded, vaithilingam2022expectation} Studies show they can support novice programmers with examples and feedback for comprehension, but they also risk encouraging shallow understanding and overreliance when generated code is accepted without reflection. \cite{kazemitabaar2023studying}

These benefits are accompanied by concerns about reliability and governance. AI-generated code may contain subtle bugs, inefficiencies, or security vulnerabilities without explicit warning, which means human review and validation remain necessary. \cite{pearce2025asleep}

\newpage
\subsection{Assistants and Agents}
Artificial intelligence systems based on large language models are commonly described in terms of two roles: AI assistants and AI agents. \cite{IBM2026AIAgentsVsAIAssistants} Although this distinction is useful for conceptual clarity, in practice the boundary between them is increasingly fluid. \cite{yao2022react} Both categories rely on similar underlying technologies, yet they differ in how these capabilities are organized, how control is distributed, and how tasks are executed within the system.

AI assistants are interactive systems that respond to user input in a direct and context-aware manner. Their primary role is to support users by generating relevant outputs, such as text or code, based on prompts. The interaction model is fundamentally reactive: each response is triggered by explicit user input, and the system does not independently pursue goals beyond the current exchange. While assistants may retain short-term conversational context, their behavior is largely constrained to a prompt-response cycle. As a result, they are well suited for tasks that require human oversight, iterative refinement, and clear user control, such as programming support, writing assistance, or question answering.

AI agents, in contrast, represent a shift toward more autonomous and goal-oriented systems. In this approach, a language model is embedded within a broader architecture that enables planning, state management, and interaction with external tools. Agents can decompose complex objectives into smaller steps, execute these steps sequentially, and adapt their behavior based on intermediate results. This is achieved through the integration of additional components, such as memory mechanisms for maintaining context across interactions, tool interfaces for accessing external functionality, and control loops that govern decision-making. Consequently, agents are capable of handling multistep workflows and operating with a degree of independence from direct user input.

Despite these conceptual differences, the line between assistants and agents is increasingly blurred. Modern systems often combine characteristics of both, resulting in hybrid architectures. An assistant may incorporate agent-like features, such as maintaining longer-term context, suggesting proactive actions, or coordinating subtasks behind the scenes. Conversely, agents may expose assistant-like interfaces when interacting with users, relying on dialogue to clarify goals or refine outputs. This convergence reflects a broader trend in AI system design, where interaction and autonomy are not treated as mutually exclusive, but rather as complementary capabilities.

\subsection{Specialized Models}
At the core of AI assistants lie foundation models based on the transformer architecture introduced by Attention Is All You Need. \cite{vaswani2017attention} Building upon these are an evolution of the general-purpose language models, often referred to as code models. These models are specialized models tailored for software engineering tasks. While they are built upon the same underlying architectures as general large language models, they are trained on large corpora of code and natural language, allowing them to learn statistical patterns that generalize across coding tasks. \cite{chen2021evaluating}

The primary benefit of this specialization is the acquisition of domain-specific knowledge. Code models are better equipped than general LLMs to understand programming constructs, adhere to syntax and semantic constraints, and generate outputs that reflect software engineering practices. This results in improved performance across a wide range of development tasks and enables their integration into tools that support developer productivity, such as intelligent code completion and automated documentation generation. \cite{chen2021evaluating ,li2022competition} However, coding models suffer from the same challenges as other LLMs, including hallucinations, outdated knowledge, and difficulties handling project-specific or proprietary contexts. \cite{ji2023survey, gao2023retrieval} These challenges have motivated the integration of external knowledge sources.

\newpage
\subsection{Retrieval-Augmented Generation} \label{sec:related-work:RAG}
Retrieval-Augmented Generation (RAG) is a technique for enhancing LLMs by enabling them to access and incorporate external knowledge at inference time. \cite{lewis2020retrieval} Rather than relying solely on the information encoded in model parameters during training, RAG systems retrieve relevant documents from a knowledge source, such as a database, document collection, or code repository, and use this information to inform the generated response. At a high level, RAG operates in two stages: a retrieval step, where a user query is used to identify relevant pieces of information, and a generation step, where the retrieved content is incorporated into the model's input context to produce a response. This combination allows RAG systems to improve factual accuracy, remain up-to-date without retraining, and ground their outputs in verifiable sources. Over time, this basic retrieve-then-generate paradigm has evolved into a range of increasingly sophisticated architectures that extend how retrieval is performed, integrated, and controlled. Recent surveys categorize this evolution into distinct categories, naive RAG, advanced RAG, modular RAG, and agentic RAG, each addressing specific limitations of earlier designs. \cite{gao2023retrieval, singh2025agentic} Naive RAG treats retrieval as a single, static preprocessing step; advanced RAG refines it through multi-stage pipelines and iterative reasoning; modular RAG generalizes the architecture into flexible components; and agentic RAG elevates retrieval to a dynamic, decision-driven process embedded within autonomous agents.

\subsection{Agentic Capabilities for AI Assistants}
As mentioned above, the line between AI assistants and AI agents are increasingly blurred. The emergence of agentic capabilities marks a significant shift in the design and functionality of AI assistants within software engineering. While early generations of assistants were predominantly reactive, limited to single-turn interactions and narrow task execution, recent advances in LLMs have enabled systems that exhibit goal-directed, iterative, and context-aware behavior. These agentic systems extend beyond prompt-response paradigms by incorporating mechanisms for planning, tool use, and self-reflection, thereby aligning more closely with traditional notions of autonomous agents in artificial intelligence. This evolution reflects a convergence between interactive assistants and autonomous software agents, where the distinction increasingly depends on the degree of autonomy and environmental interaction.

A central characteristic of agentic AI assistants is their ability to decompose complex tasks into sequences of intermediate steps and to dynamically adapt these plans based on feedback from the environment. Approaches such as ReAct integrate reasoning and acting by interleaving natural language "thoughts" with concrete actions, enabling models to query external systems, retrieve information, or execute code as part of their problem-solving process. \cite{yao2022react} This paradigm has proven particularly effective in software engineering contexts, where tasks such as debugging, test generation, and refactoring require iterative exploration and validation.

Tool augmentation is a foundational enabler of these capabilities. By granting AI assistants access to external tools, such as compilers, interpreters, static analysis frameworks, and documentation retrieval systems, researchers have significantly expanded the operational scope of LLM-based systems. In this setting, tool invocation becomes an integral part of the reasoning process, requiring the model to select appropriate tools, construct valid inputs, and interpret outputs correctly. Frameworks for tool-augmented reasoning formalize this interaction, often treating tool selection as a decision-making problem guided by the model's internal representations. \cite{schick2023toolformer} While this integration enhances practical utility, it also introduces new challenges related to robustness and security. Incorrect tool usage can lead to failures, and unrestricted access to external systems raises concerns about data leakage, privilege escalation, and unintended side effects.

Another important dimension of agentic systems is their capacity for self-reflection and iterative improvement. Techniques such as Reflexion and self-consistency prompting enable models to evaluate their own outputs and revise them in subsequent iterations, effectively implementing a form of meta-reasoning. \cite{shinn2023reflexion} In software engineering applications, this capability is particularly valuable for tasks that benefit from trial-and-error, such as synthesizing correct programs or generating test cases that satisfy specific coverage criteria. By maintaining internal feedback loops, agentic assistants can progressively refine their outputs, thereby improving reliability without requiring external supervision at every step.

The concept of multi-agent collaboration further extends the agentic paradigm by distributing responsibilities across multiple specialized agents. In such systems, distinct agents may be responsible for planning, implementation, testing, and verification, communicating through structured protocols or shared memory representations. \cite{hong2023metagpt} This modularization can improve scalability and allow for more interpretable workflows, as each agent operates within a well-defined role. However, coordination among agents introduces additional complexity, including synchronization overhead, communication bottlenecks, and the risk of error propagation across agent boundaries.


\section{Robotic Application Development}
An automation solution for an industrial task consists of both hardware and software. The hardware may include the robot manipulator, grippers, peripheral devices, and safety components, while the software spans both system software and the user program that makes the specific application perform the intended task.

Several approaches have attempted to make the robot programming easier. User-interaction-based Programming by Demonstration \cite{ravichandar2020recent} and abstraction-based Programming by Skills \cite{saukkoriipi2020programming}, aim to raise the level of abstraction in robot programming and reduce the amount of low-level engineering work required from the user. However, these lines of work also illustrate a persistent limitation: it remains difficult to formalize and capture the broad practical knowledge needed to make such approaches work reliably at scale.

In robotics, software reuse has also been explored as a way to accelerate development and improve system reliability. Much of this effort has focused on ROS, the open-source Robot Operating System. \cite{quigley2009ros} Today, the ROS ecosystem allows roboticists to assemble complete robotic systems by integrating existing software modules instead of developing all functionality from the ground up. \cite{kolak2020takes} However, this approach to reuse is primarily aimed at expert roboticists and program reuse remains difficult in practice. Estefo et al. surveyed 119 ROS users and found that more than 70\% had experienced failing to successfully reuse existing packages and participants reporting the search process as time-consuming and often unproductive, often due to challenges such as abandoned packages. \cite{estefo2019robot}

In industrial applications, Universal Robots introduced "URCaps", software extensions that simplify the integration of peripherals into robot systems. While this reduces the complexity of building automation applications, it does not directly support the reuse of previously developed application code.

Recent advances in machine learning and large-scale foundation models have nevertheless renewed interest in automating parts of robotic application development. A prominent example is SayCan, which combines large language models with robot affordance functions to generate high-level action sequences grounded in the robot's actual capabilities. \cite{brohan2023can} In SayCan, the language model proposes candidate actions while learned value functions estimate whether these actions are feasible in the current environment. This enables the system to decompose abstract natural-language instructions into executable robot behaviors. The work demonstrates that large language models can contribute semantic and procedural knowledge to robotic planning, but it also highlights a limitation: the resulting solutions still depend on predefined low-level skills and carefully engineered grounding mechanisms. The generated behavior therefore remains constrained by the available skills, training data, and environment-specific affordance models.

Subsequent work has extended this direction toward increasingly general robot policies. RT-1 introduced a large-scale transformer-based robotic control architecture trained on 130,000 real-world demonstrations covering more than 700 tasks. \cite{brohan2022rt1} Rather than focusing only on symbolic task planning, RT-1 directly maps vision and language inputs to robot actions using a unified transformer architecture. The study demonstrated improved robustness and generalization across tasks, environments, and objects, suggesting that large and diverse datasets can support more reusable robot behaviors.

RT-2 further extends this approach by integrating vision-language models trained on Internet-scale data directly into robotic control. \cite{zitkovich2023rt2} The model represents robot actions as textual tokens, allowing robotic trajectories and web-scale vision-language data to be trained within a shared framework. This enables forms of semantic reasoning and generalization that were previously difficult to achieve in robotics, including reasoning about novel objects, symbols, and relationships between objects. RT-2 demonstrates that pretrained vision-language models can transfer broad semantic knowledge into robotic behavior generation.

Together, these approaches represent a significant shift in robotics research from manually engineered task logic toward data-driven policy generation and semantic reasoning. Such systems are particularly promising for robots operating in dynamic and unpredictable environments, where flexibility, natural-language interaction, and adaptation to novel situations are important capabilities. However, industrial automation environments often prioritize different requirements, including deterministic execution, repeatability, and compatibility with established safety and validation procedures. Consequently, despite major advances in embodied AI and robotic foundation models, the development of industrial robotic applications still relies heavily on explicit robot programs and engineering expertise. This suggests that there remains a need for approaches that support the creation, comprehension, and reuse of robot program code within existing industrial development practices.


% \section{Contribution}
% Consequently, this thesis aims to contribute to the existing body of knowledge by investigating how large language models and robotic data can support the development and understanding of robotic programs.